\begin{thebibliography}{99}

\bibitem{fama1970}
Fama, E. F. (1970). Efficient capital markets: A review of theory and empirical work. \textit{Journal of Finance}, 25(2), 383--417.

\bibitem{jegadeesh1993}
Jegadeesh, N., \& Titman, S. (1993). Returns to buying winners and selling losers: Implications for stock market efficiency. \textit{Journal of Finance}, 48(1), 65--91.

\bibitem{fama1992}
Fama, E. F., \& French, K. R. (1992). The cross-section of expected stock returns. \textit{Journal of Finance}, 47(2), 427--465.

\bibitem{ang2006}
Ang, A., Hodrick, R. J., Xing, Y., \& Zhang, X. (2006). The cross-section of volatility and expected returns. \textit{Journal of Finance}, 61(1), 259--299.

\bibitem{mclean2016}
McLean, R. D., \& Pontiff, J. (2016). Does academic research destroy stock return predictability? \textit{Journal of Finance}, 71(1), 5--32.

\bibitem{evstigneev2002}
Evstigneev, I. V., Hens, T., \& Schenk-Hopp\'{e}, K. R. (2002). Market selection of financial trading strategies: Global stability. \textit{Mathematical Finance}, 12(4), 329--339.

\bibitem{evstigneev2006}
Evstigneev, I. V., Hens, T., \& Schenk-Hopp\'{e}, K. R. (2006). Evolutionary stable stock markets. \textit{Economic Theory}, 27(2), 449--468.

\bibitem{evstigneev2015}
Evstigneev, I. V., Hens, T., \& Schenk-Hopp\'{e}, K. R. (2015). \textit{Mathematical Financial Economics: A Basic Introduction}. Springer.

\bibitem{hens2005}
Hens, T., \& Schenk-Hopp\'{e}, K. R. (2005). Evolutionary finance: Introduction to the special issue. \textit{Journal of Mathematical Economics}, 41(1-2), 1--5.

\bibitem{scholl2020}
Scholl, M. P., Calinescu, A., \& Farmer, J. D. (2020). How market ecology explains market malfunction. \textit{arXiv preprint arXiv:2009.09454}.

\bibitem{scholl2022}
Scholl, M. P., et al. (2022). Towards evology: A market ecology agent-based model of US equity mutual funds I. \textit{arXiv preprint arXiv:2210.11344}.

\bibitem{scholl2023}
Scholl, M. P., et al. (2023). Towards evology: A market ecology agent-based model of US equity mutual funds II. \textit{arXiv preprint arXiv:2302.01216}.

\bibitem{finevo2026}
FinEvo Contributors. (2026). FinEvo: From isolated backtests to ecological market games for multi-agent financial strategy evolution. \textit{arXiv preprint arXiv:2602.00948}.

\bibitem{lo2004}
Lo, A. W. (2004). The adaptive markets hypothesis: Market efficiency from an evolutionary perspective. \textit{Journal of Portfolio Management}, 30(5), 15--29.

\bibitem{lo2005}
Lo, A. W. (2005). Reconciling efficient markets with behavioral finance: The adaptive markets hypothesis. \textit{Journal of Investment Consulting}, 7(2), 21--44.

\bibitem{lo2012}
Lo, A. W. (2012). Adaptive markets and the new world order. \textit{Financial Analysts Journal}, 68(2), 18--29.

\bibitem{lo2017}
Lo, A. W. (2017). \textit{Adaptive Markets: Financial Evolution at the Speed of Thought}. Princeton University Press.

\bibitem{ito2012}
Ito, M., \& Sugiyama, S. (2012). A test of the adaptive market hypothesis using a time-varying AR model in Japan. \textit{arXiv preprint arXiv:1207.1842}.

\bibitem{arthur1997}
Arthur, W. B., Holland, J. H., LeBaron, B., Palmer, R., \& Tayler, P. (1997). Asset pricing under endogenous expectations in an artificial stock market. In \textit{The Economy as an Evolving Complex System II} (pp. 15--44). Addison-Wesley.

\bibitem{lebaron2002}
LeBaron, B. (2002). Building the Santa Fe artificial stock market. \textit{Working Paper, Brandeis University}.

\bibitem{ehrentreich2007}
Ehrentreich, N. (2007). \textit{Agent-Based Modeling: The Santa Fe Institute Artificial Stock Market Model Revisited}. Springer.

\bibitem{yang2015}
Yang, L., et al. (2015). A comparison of US and Chinese financial market microstructure: Heterogeneous agent-based multi-asset artificial stock markets approach. \textit{Journal of Economic Interaction and Coordination}, 10(2), 277--305.

\bibitem{trimborn2018}
Trimborn, T., \& Frank, M. (2018). SABCEMM: A simulator for agent-based computational economic market models. \textit{arXiv preprint arXiv:1801.01811}.

\bibitem{trimborn2018stylized}
Trimborn, T., Otneim, H., \& Frank, M. (2018). Simulation of stylized facts in agent-based computational economic market models. \textit{arXiv preprint arXiv:1812.02726}.

\bibitem{kreuser2024}
Kreuser, J., \& Trimborn, T. (2024). Robust mathematical formulation of agent-based models. \textit{arXiv preprint}.

\bibitem{farmer2002}
Farmer, J. D., \\& Joshi, S. (2002). The price dynamics of common trading strategies. \textit{Journal of Economic Behavior \& Organization}, 49(2), 149--171.

\bibitem{farmer2005}
Farmer, J. D., Patelli, P., \\& Zovko, I. I. (2005). The predictive power of zero intelligence in financial markets. \textit{Proceedings of the National Academy of Sciences}, 102(6), 2254--2259.

\bibitem{brock1997}
Brock, W. A., \& Hommes, C. H. (1997). A rational route to randomness. \textit{Econometrica}, 65(5), 1059--1095.

\bibitem{brock1998}
Brock, W. A., \& Hommes, C. H. (1998). Heterogeneous beliefs and routes to chaos in a simple asset pricing model. \textit{Journal of Economic Dynamics and Control}, 22(8-9), 1235--1274.

\bibitem{lux1999}
Lux, T., \& Marchesi, M. (1999). Scaling and criticality in a stochastic multi-agent model of a financial market. \textit{Nature}, 397(6719), 498--500.

\bibitem{lux2000}
Lux, T., \& Marchesi, M. (2000). Volatility clustering in financial markets: A microsimulation of interacting agents. \textit{International Journal of Theoretical and Applied Finance}, 3(4), 675--702.

\bibitem{contr2001}
Cont, R. (2001). Empirical properties of asset returns: Stylized facts and statistical issues. \textit{Quantitative Finance}, 1(2), 223--236.

\bibitem{contr2007}
Cont, R. (2007). Volatility clustering in financial markets: Empirical facts and agent-based models. In \textit{Long Memory in Economics} (pp. 289--309). Springer.

\bibitem{chiarella2009}
Chiarella, C., Iori, G., \\& Perell\'{o}, J. (2009). The impact of heterogeneous trading rules on the limit order book and order flows. \textit{Journal of Economic Dynamics and Control}, 33(3), 525--537.

\bibitem{raberto2001}
Raberto, M., Cincotti, S., Focardi, S. M., \& Marchesi, M. (2001). Agent-based simulation of a financial market. \textit{Physica A}, 299(1-2), 319--327.

\bibitem{chen2012}
Chen, S. H., Chang, C. L., \& Du, Y. R. (2012). Agent-based economic models and econometrics. \textit{Knowledge Engineering Review}, 27(2), 187--219.

\bibitem{lebaron2006}
LeBaron, B. (2006). Agent-based computational finance. In \textit{Handbook of Computational Economics} (Vol. 2, pp. 1187--1233). Elsevier.

\bibitem{tesfatsion2006}
Tesfatsion, L., \& Judd, K. L. (Eds.). (2006). \textit{Handbook of Computational Economics, Vol. 2: Agent-Based Computational Economics}. Elsevier.

\bibitem{hommes2006}
Hommes, C. H. (2006). Heterogeneous agent models in economics and finance. In \textit{Handbook of Computational Economics} (Vol. 2, pp. 1109--1186). Elsevier.

\bibitem{kirman1993}
Kirman, A. (1993). Ants, rationality, and recruitment. \textit{Quarterly Journal of Economics}, 108(1), 137--156.

\bibitem{farmer2024}
Farmer, J. D. (2024). Market ecology and the evolution of trading strategies. \textit{Annual Review of Financial Economics}, 16.

\bibitem{kaplan2024}
Kaplan, S. N., \& Schoar, A. (2005). Private equity performance: Returns, persistence, and capital flows. \textit{Journal of Finance}, 60(4), 1791--1823.

\bibitem{berk2004}
Berk, J. B., \& Green, R. C. (2004). Mutual fund flows and performance in rational markets. \textit{Journal of Political Economy}, 112(6), 1269--1295.

\bibitem{pastor2015}
P\`{a}stor, L., Stambaugh, R. F., \& Taylor, L. A. (2015). Scale and skill in active management. \textit{Journal of Financial Economics}, 116(1), 23--45.

\bibitem{novy2022}
Novy-Marx, R., \& Velikov, M. (2022). Assaying anomalies. \textit{Working Paper}.

\bibitem{hsu2023}
Hsu, J., et al. (2023). Factor crowding and the decay of factor returns. \textit{Financial Analysts Journal}, 79(3).

\bibitem{gabaix2006}
Gabaix, X., Gopikrishnan, P., Plerou, V., \& Stanley, H. E. (2006). Institutional investors and stock market volatility. \textit{Quarterly Journal of Economics}, 121(2), 461--504.

\bibitem{stein2009}
Stein, J. C. (2009). Sophisticated investors and market efficiency. \textit{Journal of Finance}, 64(4), 1517--1548.

\bibitem{garleanu2013}
G\^{a}rleanu, N., \& Pedersen, L. H. (2013). Dynamic trading with predictable returns and transaction costs. \textit{Journal of Finance}, 68(6), 2309--2340.

\bibitem{pedersen2015}
Pedersen, L. H. (2015). \textit{Efficiently Inefficient: How Smart Money Invests and Market Prices Are Determined}. Princeton University Press.

\bibitem{akepanidtaworn2023}
Akepanidtaworn, K., Di Mascio, R., Imas, A., \& Schmidt, L. (2023). Selling fast and buying slow: Heuristics and trading performance of institutional investors. \textit{Journal of Finance}, 78(6), 3293--3342.


\bibitem{zhangwei2013}
张维, 武自强, 张永杰, 熊熊, 等. (2013). 基于复杂金融系统视角的计算实验金融：进展与展望. \textit{管理科学学报}, 16(6), 85--94.

\bibitem{zhangwei2012methodology}
张维, 李悦雷, 熊熊, 张永杰, 张小涛. (2012). 计算实验金融的思想基础与研究范式. \textit{系统工程理论与实践}, 32(3), 495--507.

\bibitem{zhangwei2009}
张维, 李根, 熊熊, 韦立坚, 王雪莹. (2009). 资产价格泡沫研究综述：基于行为金融和计算实验方法的视角. \textit{金融研究}, (8), 182--193.

\end{thebibliography}
